\documentclass[11pt]{article}
\usepackage[a4paper,margin=2.6cm]{geometry}
\usepackage{amsmath,amssymb,amsthm}
\newtheorem{lemma}{Lemma}
\theoremstyle{remark}\newtheorem*{remark}{Remark}
\title{Arithmetic SOS for the truncated Weil form:\\
\large what exists, what the smallest window shows, what would be a factor}
\author{D.\ Joubert}
\date{4 September 2026}
\begin{document}
\maketitle
\begin{abstract}
\noindent
On a window where $Q_L\succ0$, a sum-of-squares factorization exists:
$Q_L=\|Q_L^{1/2}\cdot\|^2$. That factor is spectral (or Cholesky). It is
not built from the primes. An \emph{arithmetic} SOS is a factor
$T_L$ written from $P$, $A$ and the towers $T_p$ only, with
$T_L^*T_L=Q_L$. At $\mu=3$ --- one prime, one negative mode of
$P+A$ --- the Cholesky factor on $V_9$ is dense, $T_2$ is indefinite,
and completing the square against the lag at $\log2$ does not produce
an explicit $T_L$. The closed SOS that does exist is the zero-side
Gram. An arithmetic SOS for every $L$ is point~4 of the notebook: RH.
\end{abstract}

\section{Three factorizations, only one arithmetic in intent}
On $V_N$ at a window with $Q\succ0$,
\[
Q=\sum_{j=1}^{N}\lambda_j v_jv_j^{\!\top}
=C^{\!\top}C
=\sum_{\gamma}\hat\eta(\gamma)\,\hat\eta(\gamma)^{\!\top}.
\]
The first two are $Q^{1/2}$ (spectral / Cholesky). They were computed
at $\mu=11$ ($47$ squares, certified by congruence, not by ball
Cholesky). The third is $T_{\mathrm{z}}f=(\hat f(\gamma))_{\gamma}$,
closed form, and uses the zeros. An arithmetic $T_L$ would satisfy
$T_L^*T_L=P+A-\sum T_p$ with each row of $T_L$ a local pairing
(archimedean kernel or lag $\log p^k$), no $\gamma$ and no
eigenvector of $Q$.

\section{The smallest laboratory: $\mu=3$}
Only the tower of $2$ is present. Spectro on $V_9$:
\[
\lambda_{\min}(P+A)=-0.073,\quad
\operatorname{spec}(T_2)\subset[-0.490,0.490],
\quad
\lambda_{\min}(Q)=1.03\times10^{-7}.
\]
$P+A$ has \emph{one} negative direction. $T_2$ is indefinite of
half-dimension (four negative eigenvalues on $V_9$). $Q$ is positive
with the known razor. The numerical Cholesky factor $C$ of $Q$ on
$V_9$ is dense in the lower triangle: no band, no visible Toeplitz,
no displacement rank that would write $C$ as a short combination of
the $\Theta(\log2)$ matrix and the archimedean kernel.

\paragraph{Completing the square fails to be explicit.}
Write $P+A=A_+-\lambda_-uu^{\!\top}$ with $\lambda_-=0.073$. Then
\[
Q=A_+-\lambda_-uu^{\!\top}-T_2.
\]
This is a difference of an PSD piece and an indefinite piece. It is
not $A_+$ plus a square built from $U_{\log2}$, because $T_2$ itself
is not a square: $\Theta_f(\log2)=\langle f,\tau_{\log2}f\rangle$ is a
correlation, not a norm. Expanding $\tau_{\log2}$ in the cosine basis
recovers the matrix of $T_2$, not a list of arithmetic linear forms
whose Gram is $Q$.

\section{What an arithmetic $T_L$ would look like}
A candidate with the right shape, not known to work:
\[
(T_Lf)_p
=\sqrt{\log p}\,p^{-1/4}\bigl(\alpha_p f+\beta_p\tau_{\log p}f\bigr)
\quad\text{plus an archimedean row $T_\infty f$.}
\]
Then $\|T_Lf\|^2$ produces $\Theta(\log p)$ cross terms and
$\|f\|^2$, $\|\tau f\|^2$ squares. Matching the coefficients of
$P+A-\sum T_p$ is a finite linear algebra problem on each window ---
and a moving target as $L$ grows (new primes, new lags $k\log p$).
At $\mu=3$ this ansatz has four real scalars $(\alpha_2,\beta_2)$ plus
the archimedean row. It does not have enough parameters to reproduce
a $9\times9$ PSD matrix of rank $9$. Adding internal lags
$k=2,3,\ldots$ does not help: $2^k\le3$ forbids them.

The zero-side factor has as many rows as there are $\gamma$ in band,
which is the right count. The Euler side has as many rows as there
are prime powers $\le\mu$, which is too few (four towers at
$\mu=11$, against $47$ dimensions). An arithmetic SOS must therefore
use a \emph{continuous} archimedean fibre, not only the discrete
lags. That fibre is the Weil archimedean integral, which is not
itself a square for $L>0.85$. The mixing of a non-square continuous
piece with four discrete lags is the whole problem.

\section*{Status}
Explored, not constructed. On every finite $V_N$ with $Q\succ0$ an
SOS exists and is $Q^{1/2}$. The Cholesky of the smallest PSD window
is unstructured. The Euler lags are too few to span the factor. The
closed SOS is the list of evaluators at the zeros. Passing from that
list to a factor written from $P,A,T_p$ for every $L$ is RH.
\end{document}
